\documentclass[11pt,a4paper]{article}
\usepackage[margin=2.5cm]{geometry}

\usepackage[T1]{fontenc}   
\usepackage[utf8]{inputenc} 
\usepackage[english]{babel}

\usepackage{newtxtext,newtxmath}

\usepackage{microtype}

\usepackage{xcolor}
\definecolor{citationred}{RGB}{150,0,0}

\usepackage[
  colorlinks=true,
  citecolor=citationred,
  linkcolor=black,
  urlcolor=black
]{hyperref}

\usepackage[authoryear,round]{natbib}



\usepackage[most]{tcolorbox}  
\tcbuselibrary{skins}         


\newtcolorbox{commentbox}{
  colframe=green,        
  boxrule=0.4pt,         
  sharp corners,         
  enhanced,              
  left=2mm,right=2mm,    
  bottom=2mm,
  top=1.9em,             
  overlay={
    \node[anchor=north west,fill=citationred,text=white,
          font=\bfseries,inner ysep=2pt,inner xsep=4pt]
         at (frame.north west) {Commentaires};
  },
}

\renewcommand{\thefootnote}{\fnsymbol{footnote}}

\begin{document}

\thispagestyle{empty}

\begin{center}
    {\Large\bfseries
    The Acceptability of Internationally Redistributive Policies
    }\\[0.5cm]

    {\small
    Adrien Fabre\footnotemark[1]
    \qquad
    Lubin Bénéteau % AF: ajoute ton affiliation
    }\\[0.2cm]

    {\footnotesize Draft: \textcolor{blue}{1}}
\end{center}

\footnotetext[1]{CNRS, CIRED}
% AF: dire qu'on a utilisé ChatGPT pour la relecture à l'éditeur et lui demander s'il faut le rajouter en acknowledgments

% AF! une remarque générale est que tu as souvent dévié du plan que j'avais proposé. J'ai dit que tu pouvais dévier mais j'aurais dû dire qu'il fallait me consulter avant de le faire. Maintenant, j'aimerais que tu reprennes la structure et les analyses que j'avais proposées, et que tu modifies le chapitre en ce sens. Sauf si tu veux dévier pour une raison particulière, auquel cas on en discute. Si les consignes étaient trop sybillines dis-moi, je peux détailler. Pour donner juste deux exemples: l'intro ne suit pas le plan proposé, et la 1ere section devrait commencer par l'étendue du soutien à la redistribution mondiale dans des questions générales (ISSP, UNDP)

%============================================================
%                   DÉBUT PARTIE COMMENTAIRES
%============================================================
\vfill
\begin{tcolorbox}[
    enhanced,
    breakable,
    colback=white,
    colframe=green!75!black,
    boxrule=1pt,
    arc=2mm,
    left=4mm,
    right=4mm,
    top=7mm,
    bottom=4mm,
    attach boxed title to top left={xshift=5mm,yshift=-3mm},
    boxed title style={
        colback=green!75!black,
        colframe=green!75!black,
        boxrule=0pt,
        arc=1mm,
        left=3mm,
        right=3mm,
        top=1mm,
        bottom=1mm
    },
    fonttitle=\bfseries\large\color{white},
    title={Commentaires}
]
% \input{chapter/commentaires}

\end{tcolorbox}
\vfill
%============================================================
%                   FIN PARTIE COMMENTAIRES
%============================================================

\newpage
\setcounter{page}{1}
\section*{Introduction}

Moussa Abdou is a 31-year-old man living in a rural area near Maradi, Niger, where he works as a subsistence farmer and seasonal agricultural labourer. His monetary income is irregular and highly dependent on agricultural cycles, but can be estimated at approximately USD 45 per month. Clara Meier, by contrast, is a 34-year-old woman living in Lausanne, Switzerland, where she works as a qualified nurse in a public hospital and earns a gross monthly income of approximately USD 7,200. Their situations provide a concrete illustration of the scale of contemporary global inequality: people of similar ages may face radically different material opportunities depending largely on where they are born and live.
% AF: c'est inhabituel pour une production scientifique d'avoir une partie narrative comme ça, je pense qu'il vaut mieux illustrer l'inégalité par des statistiques*. Cela dit on peut tester auprès de l'éditeur pour voir ce qu'il en pense si tu veux susciter de l'émotion chez le lecteur. 
% * Par exemple une paraphrase de: redirecting just 1\% of high-income countries' output to low-income countries (LICs) would mechanically double their national income.\footnote{The GDP per capita of high-income countries (HICs) is \href{https://data.worldbank.org/indicator/NY.GDP.PCAP.PP.CD?contextual=default&end=2024&locations=EU-ZG-XD-XM-1W-IN-US-CD-BI-LU-CN&start=2024&view=bar}{28 times} greater than that of low-income countries (LICs) at Purchasing Power Parity (PPP) and \href{https://data.worldbank.org/indicator/NY.GDP.PCAP.CD?end=2024&locations=EU-ZG-XD-XM-1W-IN-US-CD-BI-LU-CN&start=2024&view=bar}{68 times} greater in nominal terms, from World Bank 2024 data. Given that 625 million people live in one of the 25 low-income countries %defined by GDP per capita at Purchasing Power Parity (PPP) below \$1,135 per year, 
% and 1.42 billion people in high-income countries, %with a GDP per capita of over \$13,935.().
% 1\% of HICs' GDP corresponds to 60\% of LICs' GDP using PPP values and 153\% using nominal data.
% } more than enough to finance the eradication of extreme poverty \citep{sahoo_what_2025}

Such disparities raise a question of global justice. In broad terms, global justice concerns the fair and sustainable distribution of socioeconomic and environmental resources both within and across countries, subject to the ecological constraints imposed by planetary boundaries. On this view, reducing extreme poverty and narrowing large international inequalities are necessary but insufficient objectives: distributive improvements cannot be considered sustainable if they are achieved through development paths that intensify climate change and expose already vulnerable populations to further environmental harm. Global justice therefore involves not only the distribution of income and wealth, but also access to public services, economic opportunities, environmental resources, and decision-making power.

International redistribution constitutes one possible response to these inequalities. % AF: À la place, cite https://www.tandfonline.com/doi/abs/10.1080/00220388.2025.2493807 et dit que l'élimination rapide de la pauvreté (donc le respect des ODD universellemnet adoptés) ne peut se faire sans redistribution internationale, et qu'à long terme pour avoir convergence des pays en termes de PIB par habitant d'ici 2100, le Global Justice Report / Bothe et al. 2026 trouve qu'il faut des transferts de ressources et de pouvoir de l'ordre de xxx pourcents du PIB mondial entre HIC et LIC 
The scale of the underlying disparities is substantial: the richest 1\% of individuals %in high-income countries % AF: deletion
receive 
2.5 times % AF: addition
more income than the entire bottom half of the global distribution combined, 
and the richest 10\% receive more than half of the world's income % AF: addition
\citep{chancel_world_2026}.  % AF: d'où tires-tu la stat sur le top 1% des pays riches (quelle page?). Le pb avec mes changements c'est qu'on perd la dimension between countries, le pb avec ta formulation c'est qu'elle n'était pas mise en valeur. Ptet que pour corriger ça on pourrait donner la décomposition de Theil entre inégalité within et between countries, ou bien comparer le niveau de vie moyen de HIC et LIC (ce que je propose au-dessus)
Moreover, the international distribution of resources cannot be understood solely through differences in current income or gross domestic product. Large flows of value continue to move from the Global South toward the Global North through unequal economic relations, debt servicing, tax-base erosion, and other forms of extraction \citep{hickel_plunder_2021}. Debt distress and austerity can constrain the policy choices available to low-income countries and reproduce asymmetric economic relationships \citep{fischer_return_2023}. The literature also emphasizes the historical origins of contemporary distributive obligations. Colonial institutions contributed to the formation of present economic structures and continue to influence development trajectories \citep{bhambra_colonial_2021}, while slavery and colonial extraction provide distinct grounds for reparative claims \citep{solbiac_britains_2015,darity_cumulative_2022}. Climate change adds a further dimension: because high-income countries have used a disproportionate share of the global carbon budget, they may bear particular responsibilities toward countries that have contributed less to cumulative emissions but remain highly exposed to their consequences \citep{fanning_compensation_2023}. % AF: add alongside fanning_compensation_2023 Hickel & Zoomkawala, National responsibility for exceeding the Paris Agreement temperature boundaries, Ecological Economics (2026)  (j'ai été rapporteur de cet article, il devrait paraître bientôt)
Taken together, contemporary inequality, historical responsibility, asymmetric economic relations, and climate liability provide several distinct arguments for transferring resources across national borders.

International redistribution should not, however, be conflated with all cross-border resource flows. Migrants' remittances, for example, are private household transfers rather than collectively determined redistributive policies. Officially recorded remittances to low- and middle-income countries reached approximately USD 685 billion in 2024, exceeding foreign direct investment and Official Development Assistance (ODA) combined \citep{ratha_remittances_2024}. Corporate profit shifting represents a different phenomenon: rather than a deliberate transfer toward poorer countries, it erodes national tax bases and limits governments' ability to finance public expenditure. Around 36\% of multinational profits are estimated to be shifted to tax havens, while associated global corporate-tax revenue losses have been estimated at approximately USD 500--650 billion annually, although these estimates remain methodologically contested \citep{torslov_missing_2023, cobham_global_2018}. ODA is more directly redistributive and more fully institutionalized, but substantially smaller in aggregate terms. Development Assistance Committee donors provided USD 214.6 billion in ODA in 2024 on a grant-equivalent basis, followed by a preliminary decline to USD 174.3 billion in 2025 \citep{oecd_final_oda_2025, oecd_preliminary_oda_2026}. % AF: better to express all figures as % of global GDP
These comparisons indicate that the politics of international redistribution cannot be reduced to foreign aid alone. It also concerns international taxation, fiscal sovereignty, debt, climate finance, and the institutional rules governing cross-border resource allocation.

This chapter does not seek to resolve the normative question of how much international redistribution justice ultimately requires. % AF: on this you can cite Kopczuk et al 2005
It addresses a narrower positive question: to what extent would %are 
citizens % willing to % AF: changed wording to avoid WTP
accept the policies and institutions through which international redistribution could be implemented? The available literature primarily examines respondents in high-income or net-contributing countries, although several studies provide broader cross-national or globally representative evidence. The chapter therefore focuses principally on the attitudes of citizens who are likely to bear %, or believe that they will bear, part of 
the %financial and political % AF: deletion
costs of international redistribution.

% AF: indeed je pense qu'il faut supprimer ce paragraphe
Several distinct concepts must be kept separate. \emph{Support} or \emph{approval} refers to a favorable stated attitude toward a policy. \emph{Policy acceptance} generally concerns approval of a more fully specified collective rule, including its financing and institutional design. \emph{Willingness to contribute} or \emph{willingness to pay} refers to the personal sacrifice respondents declare themselves prepared to bear, which should not be equated with an observed payment. \emph{Political salience} concerns the importance that citizens attach to an issue and the weight it receives in electoral decisions, while \emph{political prominence} refers to its visibility in party programmes, electoral campaigns, and public debate. A policy may therefore receive majority support when respondents are asked about it without becoming sufficiently important to influence their vote or generate sustained pressure on political representatives.

The chapter starts from precisely this puzzle. International redistribution appears normatively important and economically relevant, yet 
%the policies that would give it institutional form remain limited and politically under-salient. 
it remains elusive. % AF: changed wording
The argument developed here is not that citizens display unconditional enthusiasm for transferring resources abroad. % AF: clumsy to start saying what the argument is *not* (rather than what it is), change this
Rather, the literature points to a pattern of \emph{conditional acceptance}. Support tends to be greater when policies are perceived as fair, effective, fiscally credible, progressively financed, and embedded in plausible arrangements for international burden-sharing. Conversely, approval may remain politically weak when citizens doubt that resources will reach intended beneficiaries, 
perceive foreign expenditure as competing with domestic redistribution, % AF: what do you have in mind behind this?
underestimate the support of others, % AF! what do you have in mind behind this? Is there evidence of bandwagon effect?
or assign little electoral importance to the issue. % AF: I don't think low salience reduces acceptance
The chapter develops % AF: change wording: In this chapter, we develop... We... 
this argument in seven steps. It first examines moral universalism as a psychological and political foundation for extending distributive concern beyond national borders. It then turns to foreign aid, the most familiar public instrument of international redistribution, and studies 
the level of support and % AF: added
how support depends on information, perceived effectiveness, institutional credibility, ideology, and competition with domestic expenditure. The third section considers % more AF: deleted
rule-based forms of redistribution, particularly international taxation of 
rich % high-income and high-wealth % AF:
individuals funding %when revenues are allocated to 
poverty reduction and public services in low-income countries.
The analysis then moves from instruments to institutions. The fourth section examines whether citizens accept the forms of global governance required to coordinate contributions, constrain free-riding, and implement policies across countries. 
It shows that support is generally stronger for democratic and functionally limited institutions than for vague proposals involving an unspecified transfer of sovereignty. % AF: j'enlèverais cette phrase car pour les autres parties, tu n'annonces pas ce que tu trouves dans l'intro. Mais peut-être qu'il faudrait faire l'inverse : résumer systématiquement les résultats dans l'intro
The fifth section turns to international carbon pricing, %climate policy, 
which would make ordinary citizens in high-income countries bear the costs, and not only the richest people. % AF: 
% provides a demanding test because household costs, distributional consequences, and participation by other countries can be made explicit. 
The sixth section distinguishes collective policy approval from stated willingness to bear a personal contribution, paying particular attention to the gap between respondents' own willingness and their beliefs about the willingness of others. % AF: s'agit-il vraiment de WTP ou bien de soutien à une mesure? Ne pas parler de WTP s'il s'agit du second
Finally, the chapter considers why internationally redistributive and climate policies remain politically marginal despite their conditional acceptability. It examines low issue salience, factual misperceptions, second-order beliefs, nationally bounded reference groups, % AF: comprends pas ce que "nationally bounded reference groups" veut dire: nationalisme ?
fiscal competition, biased perceptions among political representatives, and weaknesses in political supply.

\section{Deep values: more people share some degree of universalism} % AF! le titre ne veut rien dire. De façon générale, changer tous les titres de sections

International redistribution relies on countries' % AF: les pays n'ont pas de volonté, je changerais cette expression. Aussi, je vois pas pourquoi parler de GNI plutôt que d'income tout simplement
willingness to transfer a share of their Gross National Income (GNI) to other countries. Such transfers may be either unilateral or multilateral. Their adoption depends on domestic institutions and elected policymakers. They therefore require sufficient public and legislative support to be authorized and funded through national budgetary processes. The chapter accordingly examines attitudes toward international redistribution across the political process, from citizens' preferences to the beliefs and choices of political and economic elites.

Economists have long recognized that group membership and identity shape economic behavior. In economic models, universalism captures the degree to which altruism is invariant to the identity or group membership of the recipient \citep{Tabellini_2008}. Universalism therefore concerns not an individual's overall level of altruism or trust, but the extent to which these dispositions decline with social, national, or geographic distance \citep{cappelen_universalism_2025,enke_moral_2023}. This distinction is central to redistribution: citizens may prioritize compatriots through domestic transfers or extend their moral circle\footnote{The notion of a moral circle refers to the set of beings or groups included within an individual's sphere of moral concern. It can range from the self and close relatives to compatriots, humanity as a whole, and, in broader versions, animals and ecosystems.} to socially and geographically distant beneficiaries abroad \citep{cappelen_universalism_2025,enke_moral_2023}.

% AF: we can remove this paragrpah
The contrast between Moussa and Clara gives concrete content to this distinction. Universalism concerns the extent to which Clara incorporates Moussa's welfare into her moral evaluation relative to the welfare of her compatriots. More particularistic preferences assign greater moral weight to national or socially proximate beneficiaries, whereas more universalistic preferences make distributive concern less sensitive to nationality and social distance.

Economists measure individual variation in altruism and trust using nationally representative and cross-national surveys, often combined with structured allocation tasks \citep{cappelen_universalism_2025,enke_moral_2023}. This section examines whether universalist preferences provide a foundation for support for international redistribution. 
The central question is whether these preferences translate into support for concrete redistribution across national borders. % AF: ? pour moi la question ici c'est surtout le degré d'universalisme
Universalism varies systematically across individuals and correlates with political orientation and socioeconomic characteristics, although these relationships differ across societies \citep{cappelen_universalism_2025,enke_moral_2023}. % AF: phrase non-informative, soit la supprimer soit la renforcer par des chiffres
The translation of universalist preferences into support for international redistribution may nevertheless be mediated by two distinct filters. First, citizens may misperceive both their relative position in the global income distribution and the amount of foreign aid already provided. Second, the relationship between universalism and policy support may depend on the broader moral and ideological structure of citizens' preferences.

A first filter concerns information and the reference groups against which individuals evaluate their economic position. Using a representative German survey experiment, \citet{fehr_your_2019} % AF! fais attention il y a plusieurs papiers (dont celui-ci) qui sont cités comme working papers au lieu de citer la version publiée
examine whether respondents misperceive their national and global income ranks and whether correcting these beliefs affects support for policies addressing global inequality. Germans underestimate their global income rank by an average of approximately 15 percentage points, while preferences for global redistribution correlate with preferences over domestic redistribution, immigration, globalization, and giving to the global poor. Correcting perceived national rank changes redistributive preferences among left-leaning respondents, % AF! it changes global or national redistributive preferences? What matters before these effects is the degree of support (and degree of universalism). Say that in Fehr et al., 90% of Germans support global redistribution > 0
whereas correcting perceived global rank has no detectable effect on support for global redistribution. 
By contrast, providing Americans with information about their relative global affluence increases both foreign-aid support and donations to international charities \citep{nair_misperceptions_2018}. % AF! by how much? it's not enough to say that an effect is significant, we need to know its magnitude. E.g.:
% Aid support in US driven by information on global distribution, because people underestimate their rank by 27 centiles and overestimate global median income by a factor 10
Taken together, the studies show that information can activate global distributive concern in some settings, while national income comparisons remain an important reference point. % AF: le dernier point (national comparison) est justifié par quoi ? il vient comme un cheveu sur la soupe

% AF: c'est plutôt ce papier qu'il faudrait citer: Moral Universalism: Measurement and Economic Relevance; en citant aussi la critique par Boyer-Kassem
A second filter concerns the moral structure that organizes political preferences across policy domains. Across Western countries, \citet{enke_moral_2023} examine why support for redistribution, health care, environmental protection, affirmative action, and foreign aid forms one policy cluster, while support for military, policing, and border control forms another. They argue that heterogeneity in moral universalism helps account for this structure. More universalist respondents prefer higher spending on welfare, environmental protection, affirmative action, and foreign aid, but lower spending on the military, policing, and border control. % AF! je comprends pas, la première phrase dit "ils montrent pourquoi on a cette corrélation", et la deuxième phrase dit juste : "il y a cette corrélation". Ça ne va pas. Il faut dire par exemple que l'universalisme est un meilleur prédicteur des opinions politiques que l'axe gauche-droite ou autre, mais en l'état les phrases n'ont pas bcp de sens. Voici des idées de quoi dire sur ce papier:
% Measures universalism as the money allocated to two equally rich in-group vs. out-group persosn. Domestic universalism: family vs. compatriot; foreign.
% “Universalism is quantitatively substantially more important for explaining policy views and ideological constraint than traditional political economy variables such as income, wealth, equityefficiency preferences or beliefs about government efficiency.” Figure 8
% Ideological constraint is intracorrelation of people’s policy views. “Ideological constraint would decrease by 25% if there was no heterogeneity in universalism. To put this magnitude into perspective, our estimations suggest that ideological constraint would not change at all if there was no heterogeneity in, for example, income, wealth, education, equity-efficiency preferences or beliefs about government efficiency.”
A shift from complete in-group favoritism to equal treatment is associated with differences in policy preferences ranging from approximately 10\% to 80\% of a standard deviation.

These moral divisions are also reflected in political competition. \citet{enke_moral_2020} examines whether electoral support reflects congruence between voters' moral values and candidates' moral rhetoric. Both voters and politicians vary in their emphasis on universalist and communal values, and candidates obtain larger vote shares where their moral appeal more closely matches local values. Since the 1960s, Democratic and Republican moral rhetoric has increasingly diverged, with Democratic politicians placing progressively greater relative emphasis on universalist concepts. Moral universalism is therefore associated not only with policy preferences but also with political rhetoric, voter alignment, and electoral outcomes. % AF: je pense qu'on peut supprimer ce paragraphe. La façon dont ils définissent "universalism" dans ce papier est très éloignée de ce qu'on a en tête dans ce chapitre. Mais si tu veux citer cet article, voici quoi dire dessus: Comment ils définissent universalism; Universalism associated with Dems; has decreased over 2008-2018; Trump most communal (least universalist) candidate ever

% Voici comment je citerais ce papier (une phrase suffit) : Enke et al. 23 measure universalism at the U.S. district level using donation data, and find that a district's universalism predicts electoral outcomes better than its income or education level
Finally, universalism is reflected not only in stated preferences but also in observed patterns of giving and redistribution. \citet{enke_universalism_2024} provide field evidence linking universalist values to charitable giving and political representation. They measure district-level universalism by estimating how sharply charitable donations decline with geographic and social distance. District universalism is strongly associated with Democratic vote shares in US House elections: a one-standard-deviation increase % AF: increase in what? define district universalism
corresponds to an approximately 13-percentage-point higher Democratic vote share. Donations from more particularistic districts are concentrated in places connected by geographic proximity or stronger social ties, whereas giving from more universalistic districts is less sensitive to either form of distance. % AF: unclear. What is particularistic?
Complementary international experimental evidence shows that distributive choices respond to need, entitlement, and self-interest, while nationality has no independent effect on allocations \citep{cappelen_needs_2013}. % AF: again, unclear. Here is a better summary: lab experiment matching North Europeans with East Africans in a dictator game: the participants acted as moral cosmopolitans and did not assign importance to nationality in their distributive choices

% AF: tu peux supprimer ce paragraphe, il n'apporte rien (et c'est maladroit de citer l'APD avant la section qui porte dessus)
Taken together, these findings indicate that universalist values make concern for distant beneficiaries psychologically and politically possible, but do not automatically translate into support for international redistribution. Their policy relevance depends on whether this moral disposition can be translated into salient and institutionally credible redistributive instruments. Recent evidence indicates that global redistributive and climate policies can attract majority support in high-income countries when their design, expected effectiveness, and burden-sharing arrangements are perceived as credible \citep{fabre_majority_2025}. Foreign aid provides the natural first test: it is the most familiar budgetary form of international redistribution, yet citizens systematically misperceive its scale \citep{kaufmann_foreign_2019}, while support depends on perceived effectiveness and burden-sharing \citep{gampfer_obtaining_2014,fabre_majority_2025}.


\section{Foreign aid: majority support for increased foreign aid if it is well spent}

% AF: ce paragraphe est inutile
Foreign aid is among the most visible and institutionalized forms of international redistribution through which donor countries seek to alleviate poverty and support development abroad. It includes Official Development Assistance (ODA) and related forms of public development finance. Its scale and allocation are largely determined through national legislative and budgetary processes. Sustained increases in foreign aid therefore require sufficient public and legislative support, making donor-country citizens' preferences politically consequential.


Universalism may provide a moral foundation for international redistribution, but policy support can remain limited when citizens are poorly informed or perceive aid as competing with domestic expenditure. % AF: encore une phrase qui n'est pas justifiée par des chiffres ou une citation, à éviter
Information about global inequality may strengthen aid support, while perceived budgetary trade-offs with domestic redistribution may weaken it. International redistribution also presents a collective-action problem: opportunities for countries to free-ride may reduce domestic support when contributions are not internationally coordinated \citep{olson_1965}. % AF: je connais pas ce papier et la phrase ne me permet pas de comprendre ce qu'il faut: est-ce que c'est un papier de théorie qui donne des mécanismes possibles dans l'abstrat ? Quand on cite un papier, il vaut mieux s'en tenir à son contenu: méthodo, résultats.
The literature accordingly identifies moral concern, information, perceived provision, institutional effectiveness, and burden-sharing as distinct determinants of support for cross-border transfers \citep{cappelen_needs_2013,nair_misperceptions_2018,kaufmann_foreign_2019,gampfer_obtaining_2014}. % AF: avant de donner les déterminants du soutien, il faut commencer par donner l'ampleur du soutien (cf. le plan proposé)

The first empirical question is whether national borders constrain distributive concern when potential beneficiaries live abroad and face greater material need. In a real-effort dictator experiment spanning Norway, Germany, Uganda, and Tanzania, \citet{cappelen_needs_2013} find that participants respond to need and entitlement % AF! c'est quoi entitlement? pas clair. Aussi, papier déjà cité au-dessus, faut supprimer l'autre citation
but assign no independent moral weight to nationality. Holding productivity approximately constant, participants from high-income countries allocated 42.7\% to low-income recipients, compared with 36.7\% to high-income recipients. These allocation choices do not, however, establish that citizens support increasing tax-financed foreign-aid budgets; the translation of cross-border concern into policy support depends partly on information. % AF: les deux parties de la phrase ne sont pas logiquement reliées: c'est par parce que l'information influe sur le soutien que celui-ci est majoritaire ou minoritaire. Quel soutien à tax-financed foreign-aid budgets trouvent-ils? Quelle information apportent-ils et qu'est-ce que ça change au soutien ?
In the United States, respondents underestimate their global income rank by 27 percentage points and overestimate median global income by a factor of ten \citep{nair_misperceptions_2018}. % AF: pourquoi le citer plus haut? cette citation me semble mieux
Providing corrective information raises support for increasing foreign aid from 12\% to 22\% and increases donations to foreign charities. % AF: Citer Nair après Kaufmann, dire que l'états-unien médian souhaite maintenir l'aide à son niveau actuel, citer la proportion qui souhaite que l'aide baisse (en plus de hausse), expliquer que cette préférence pour une baisse chez Nair s'explique par une particularité états-unienne et par le fait qu'une majorité soutient une hausse de l'aide conditionnelle
Across nearly 90,000 respondents in 24 donor countries, perceived aid averages 2.65\% of GNI despite actual provision of only 0.37\%; respondents nevertheless prefer aid above its actual level \citep{kaufmann_foreign_2019}. % AF: you mean above its perceived lebel? Pcq y a rien d'étonnant à ce qu'ils préfèrent above actual (je comprends pas le "nevertheless")
% Ce papier il a plein de résultats donc on doit le citer plus en détail, pourquoi pas reproduire leur table principale: 
% Very interesting and useful (above all Table 1 in Appendix). Shows the level of perceived and desired aid in 26 countries between 2005 and 2008. 

% Provide evidence of lobbying by richest against foreign aid (p 19). Indeed, more equal country provide more aid, but because the discrepancy with desired aid is lower, not because people are more generous. This is because of fewer lobbying by richest. (p 4)

% Shows that richer want less aid (p 3), but more unequal countries want more.

% Perceived aid explains 15% of variance of desired aid. Error (desired/perceived) is negatively correlated with income.

% The response is interpreted as ODA although the formulations are (no randomization):  
% (i) What share of national income does your country actually give in foreign aid to help development / poverty alleviation in other countries?  
% (ii) What percentage in your view should it give?

% 11 categories ranging in increasing intervals

% from 0 to higher than 25%: threshold at 0.05; 0.15; 0.35; 0.75; 1.5; 2.5; 4; 7.5; 17.5; 25.

% In most countries (incl. UK, DE, FR, ES but not U.S.) desired aid is larger than perceived. In most countries the gap between the two is small, except in the U.S. where perceived is 7.5% of GDP and preferred is 3%.Use WVS and Gallup (like Chong & Gradstein, Paxton & Knack) but have more waves and the others don't use the question on perceived aid. Shows that richer want less aid ("those in the top income quintile favour ODA (as a share of GNI) that is 0.13 percentage points lower than the preferred share for individuals in the bottom 40\% of the income distribution" after controling for perceived aid).

% Argue that this is due to political influence efforts/possibilities of the rich, as they prefer less aid (support this by a theoretical model + correlations between level of lobbying and actual aid level, controling for desired aid).

Information alone does not determine aid support; institutional design also matters. Evidence from North--South climate finance illustrates these effects particularly clearly. \citet{gampfer_obtaining_2014} find that joint allocation and broader donor burden-sharing increase support in both Germany and the United States, whereas recipient-government efficiency affects support only among US respondents. % AF: pas clair, expliquer plus clairement ce que fait le papier, e.g. Using a conjoint analysis in U.S. and Germany, show xx p.p. stronger support for mitigation/adaptation funding in LIC when cost is shared with other countries (and hence lower), also find preference for recipient countries with efficient administration, and multilateral decision-making => mais il s'agit pas d'aide au dév, si ? changer de section a priori
Compared with recipient-only decision-making, joint allocation raises proposal-selection probabilities by approximately 6 percentage points in the United States and 9 percentage points in Germany. Among US respondents, increasing other donors' contribution from 10\% to 90\% raises the probability that a proposal is selected by 27.9 percentage points. 
Consistent with these findings, 62\% % AF: of whom? say which countries
favor increasing foreign aid, but most condition that increase on aid reaching intended beneficiaries, avoiding diversion, and eliciting greater contributions from other high-income countries \citep{fabre_majority_2025}. % AF: give more detail on these results

Domestic fiscal competition imposes an additional constraint on this conditional support. Economically vulnerable left-leaning respondents remain broadly committed to aid as a moral obligation but are less supportive of increasing public expenditure, a pattern consistent with domestic budgetary trade-offs \citep{cho_ideology_2024}. % AF: rather cite by saying: Eurobarometer shows strong support that foreign aid is a moral obligation and majority support to comply with our promise to increase aid (also questions on donation).
Higher-income respondents prefer lower ODA, while public aid is associated with private giving and actual ODA is negatively correlated with lobbying and private political influence \citep{kaufmann_foreign_2019}. The resulting gap between preferred and actual aid constitutes a democratic aid deficit: public support exists, but it is conditional, fiscally contested, and filtered through unequal political influence.

Foreign aid illustrates the chapter's broader argument: citizens in high-income countries may accept international redistribution, but support depends on credible institutional safeguards and remains vulnerable to domestic fiscal competition. The analysis therefore turns to whether more rule-bound, tax-based instruments can attract equally broad or stronger support.



\section{International taxation: majority support for taxing the rich to finance LICs}

%The previous section showed that foreign aid can attract majority support when citizens perceive it as effective. %, well targeted, and supported by equitable burden-sharing. AF: deletion
The analysis therefore turns to whether citizens in high-income countries accept %more demanding forms of global redistribution, particularly 
international taxes whose burden would fall predominantly on the richest individuals. % with high income or substantial wealth. AF: changed
The survey evidence presented by \citet{fabre_public_2026} provides a direct test of public acceptance of these instruments. The study combines nationally representative samples of more than 12,000 adults across eleven high-income countries -- the United States, Japan, Russia, Saudi Arabia, France, Germany, Italy, Spain, the United Kingdom, Poland, and Switzerland -- with several survey experiments. % It examines the level and robustness of support across policy designs and country-coverage scenarios, while investigating why substantial stated support has not translated into greater political prominence. % AF: deletion

% AF: ce paragraphe n'a pas sa place ici, c'est pour la dernière section
The gap between stated support and political prominence connects the evidence on global taxation to the political-science literature on issue salience. Voters assign unequal attention and electoral weight to different policy issues. Electoral choices are instead shaped by the issues that become salient and by the positions candidates adopt on those issues \citep{repass_issue_1971,krosnick_role_1988,edwards_explaining_1995,dennison_review_2019}. Consequently, stated support for global redistribution may generate little electoral pressure unless political competition makes the issue sufficiently salient.

\paragraph{International taxation of top incomes}
\citet{fabre_public_2026} first examines support for taxing top incomes globally to finance poverty reduction in low-income countries. Respondents evaluate % AF: ils évaluent une des eux choisie aléatoirement
two variants that differ in the size of the taxable group and the progressivity of the rate schedule. The first imposes an additional 15\% tax on annual after-tax income above \$120,000 at purchasing power parity, thereby targeting the global top 1\%. The second extends taxation to the global top 3\% and applies marginal rates of 15\%, 30\%, and 45\% above the respective thresholds of \$80,000, \$120,000, and \$1 million. The variants are calibrated to close poverty gaps relative to monthly income floors of \$250 and \$400, respectively, and would generate international transfers equivalent to approximately 2\% and 5\% of nominal world income.

%Even these relatively ambitious proposals attract either absolute-majority support or support that substantially exceeds opposition. % AF: deletion
Support reaches 56\% for the top-1\% tax and 50\% for the top-3\% tax, compared with opposition rates of 25\% and 28\%, respectively. Although support falls as the tax base expands, both variants attract substantially more support than opposition, indicating that respondents do not generally reject large-scale redistribution toward low-income countries \citep{fabre_public_2026}. % Say there is majority support among affected respondents overall, though not in all countries for the top 3% variant (Figure S41)

% AF: maybe skip the Zucman tax, this it's not internationally redistributive (talk about it in the paragraph about plausible policies)
The study also examines several internationally redistributive wealth taxes, including a 2\% tax on wealth above \$1 million and a distinct billionaire wealth tax corresponding %more closely 
to the proposal of %associated with  % AF: deletion
Gabriel Zucman. % cite his G20 report
Under the high-income-country coverage scenario, % AF: this coverage is not defined, rephrase
70\% accept a 2\% tax on wealth above \$1 million that allocates 30\% of its revenue to public services in low-income countries. Relative to this scenario, global coverage raises acceptance by 5 % 4.8  AF: changed
percentage points, while narrow %er 
international coverage of just a few progressive countries 
lowers it by a statistically insignificant 1 %1.4 p
ercentage point. These comparisons indicate that acceptance is only slightly affected by %not highly sensitive to  AF: change
incomplete country participation. %Acceptance is highest for the 2\% billionaire wealth tax, at 81\%, compared with 69\% and 64\% for taxes on the global top 1\% and top 3\%, respectively \citep{fabre_public_2026}. % AF: deletion

% AF: ça doit aller en section 8, et il faut aussi mentionner les conjoint experiments du papier NHB
The study further examines whether global redistribution affects choices between hypothetical political programs. In a conjoint experiment, an international millionaire tax allocating 30\% of its revenue to healthcare and education in low-income countries increases the likelihood that a program is preferred %a program's probability of selection AF: changed
by 5 percentage points. By contrast, proposing cuts to development aid reduces it %a program's probability of selection 
by 3 percentage points. Global redistribution thus affects stated program choice, despite its limited prominence in actual electoral competition \citep{fabre_public_2026}. 

% AF: ça doit aller dans section 7
The paper further examines whether respondents are willing to allocate their own tax payments to global objectives. Overall, 41\% agree that their taxes should contribute to solving global problems, compared with 28\% who disagree. %; agreement reaches 60\% among respondents expressing a view. AF: deletion
If the party to which they feel closest joined a global movement combining climate action, millionaire taxation, and poverty reduction in low-income countries, 36\% report that they would become more likely to vote for it, compared with 17\% who would become less likely. Among non-voters who identify with a left-wing party, 46\% report that the movement would make them more likely to vote, whereas 10\% report the opposite. Although hypothetical, these responses indicate that global redistributive taxation may mobilize citizens already ideologically aligned with redistributive and climate-oriented platforms \citep{fabre_public_2026}.

% AF: il s'agit de "Plausible policies" et ça doit venir en début de section
\paragraph{Foreign aid, tax cooperation, and climate-related packages}
Substantial acceptance is not confined to taxes on personal income and wealth. Acceptance reaches 70\% for both debt relief for vulnerable countries and the % a foreign-aid 
commitment 
to allocate % equal to 
0.7\% of GDP % AF: changed
to foreign aid. Acceptance rises to 79\% for the Bridgetown Initiative, %presented as; cite it
a structural reform of international development and climate 
finance. A proposal to raise the global minimum corporate-tax rate from 15\% to 35\% receives 68\% acceptance. Public acceptance therefore encompasses not only explicit redistributive instruments but also institutional reforms of development finance and international tax cooperation \citep{fabre_public_2026}.

% ça doit aller dans section 7
Support also remains substantial when international redistribution forms part of a broader climate-policy package: 68\% prefer a cooperative scenario combining millionaire taxation, vehicle electrification, doubling prices of heating gas, beef and flights, %higher prices for carbon-intensive consumption,  AF: changed
and warming limited to +2\textdegree{}C over a status quo leading to +3\textdegree{}C warming by 2100 \citep{fabre_public_2026}. The result demonstrates substantial acceptance of a costly package jointly addressing climate mitigation, international cooperation, and global redistribution, although it does not identify which component drives support.

% Overall, \citet{fabre_public_2026} finds substantial support for international taxes on top incomes and wealth when revenues are allocated to poverty reduction, public services, or climate action in low-income countries. AF: changed
Overall, \citet{fabre_public_2026} finds substantial support for international transfers when financed by  taxes on top incomes and wealth. 
% AF: pas d'accord avec cette conclusion: taxer les riches pour réduire la pauvreté à l'échelle mondiale, c'est précisément ça le redistribution mondiale. C'est les résultats sur foreign aid qui montre un soutien conditionnel (à l'absence de corruption)
% The findings nevertheless fall short of demonstrating broad and unconditional enthusiasm for global redistribution. The evidence instead indicates conditional acceptance, with support concentrated on taxes targeting very wealthy households and linked to international cooperation, poverty reduction, and climate action. 
% The central puzzle is thus the persistence of political under-provision and low issue salience despite majority or plurality support for several international redistributive policies.

\section{Global governance: majority support for global democracy on world issues}

The preceding sections examined whether citizens accept specific forms of international redistribution. A distinct question is whether they also accept the supranational institutions needed to coordinate national contributions, limit free-riding, and implement common policies across countries.

% AF: useless
The evidence reviewed thus far indicates that citizens in surveyed countries may accept international redistribution through foreign aid, taxation, and other cross-border transfers. Public acceptance is only a first step, however, because these policies also require institutional arrangements capable of coordinating contributions and ensuring effective implementation. 
As \citet{kopczuk_limitations_2005} argue, decentralized national redistributive systems have limited effects on global inequality, while voluntary international transfers remain constrained by free-riding. Some degree of global coordination may therefore be necessary to reduce extreme poverty and address inequalities extending beyond national borders. This section examines whether citizens accept the forms of global governance required to overcome these coordination problems.

% AF: je vois pas l'intérêt du paragraphe suivant non plus, ça ne parle pas des attitudes
The need for coordination extends beyond international redistribution. Global inequality, climate change, pandemics, armed conflict, and migration create cross-border effects that national policies cannot address effectively in isolation. These challenges have generated an extensive literature on multilateral agreements, international organizations, and institutional reform. Global governance, however, need not entail a unitary world state. It can instead be understood as the rules, institutions, organizations, and decision-making procedures through which countries manage genuinely transnational problems. Under this conception, global governance does not replace domestic policymaking but addresses problems whose externalities and coordination requirements justify action above the national level \citep{ghassim_who_2024,dryzek_global_2011}.

% AF: pareil, ça parle pas d'attitudes
Public acceptance should therefore depend on the institutional design and scope of supranational authority. In particular, citizens may distinguish between an unspecified transfer of sovereignty and an institution that is democratically organized and restricted to problems requiring international coordination. \citet{dryzek_global_2011} provide a deliberative framework for evaluating such legitimacy, emphasizing the participation of affected actors, deliberation across competing perspectives, and effective transmission between civil society and authoritative decision-making. Global institutions must consequently be assessed not only by the geographical level at which they operate, but also by the democratic procedures through which authority is exercised. Their establishment ultimately depends on whether citizens accept both their mandate and their institutional design.

\citet{ghassim_who_2024} examine attitudes toward world government using survey data from approximately 42,000 respondents across 17 countries collected between 2017 and 2021. % AF: say which countries (perhaps in footnote)
Their experimental design varies two dimensions of legitimacy. Institutional legitimacy concerns whether the proposed global authority is explicitly democratic, whereas functional legitimacy concerns whether its mandate is confined to problems requiring global coordination. The results show that both dimensions substantially shape % AF: increase?
public acceptance.

With equal country weights, 52\% of respondents reject an unspecified world government, whereas population weighting yields 58\% support. When the proposed government is explicitly democratic, support rises to 67\% with equal country weights and 71\% with population weights. Restricting its mandate to global issues produces corresponding support rates of 66\% and 69\%, while a government focused specifically on COVID-19 receives 64\% and 69\% support. Combining institutional and functional legitimacy produces still higher acceptance: a democratic world government focused on global issues receives 69\% support with equal country weights and 72\% with population weights, while a democratic government focused on COVID-19 receives 72\% and 74\%, respectively \citep{ghassim_who_2024}.

The democratic global-issues proposal attracts majority support in every surveyed country except the United States. Support ranges from 75\% to 82\% in Egypt, India, Kenya, Indonesia, South Korea, Colombia, and Hungary; apart from the United States, the lowest levels are recorded in Russia, at 56\%, and Argentina, at 58\%. Citizens therefore respond not only to the transfer of authority beyond the nation-state but also to whether that authority is democratically constituted and functionally confined to genuinely global problems \citep{ghassim_who_2024}.

% AF: tu peux retirer ce paragraphe
The evidence thus traces a progressive expansion in the domain of acceptable redistribution. Universalistic values make concern for distant populations possible; attitudes toward foreign aid and international taxation show that this concern can extend to concrete transfers; and the findings of \citet{ghassim_who_2024} indicate that it may also extend to the institutions needed to govern transnational problems. The next question is whether this acceptance persists when respondents evaluate specific institutional reforms rather than an abstract model of global authority.

\citet{fabre_public_2026} examines whether acceptance of global governance extends beyond an abstract world government to concrete forms of international cooperation and institutional reform. When asked whether governments should cooperate to promote convergence in GDP per capita across countries by the end of the century, 61\% of respondents answer positively and 26\% negatively. Excluding non-responses, the overall acceptance rate reaches 70\%, although it is lower in the United States, at 56\%. Reform of the United Nations Security Council % AF: détailler quelle réforme
likewise receives 69\% acceptance, indicating that citizens may support not only internationally redistributive outcomes but also changes to the institutional architecture through which global decisions are made.

% AF: déjà dit plus haut, enlever
This acceptance may also have electoral implications. If the party to which they feel closest joined a global movement promoting climate action, millionaire taxation, and poverty reduction in low-income countries, 36\% of respondents report that they would become more likely to vote for it, compared with 17\% who would become less likely. Among non-voters who identify with a left-wing party, the corresponding figures are 46\% and 10\%. Although these responses concern hypothetical voting intentions, they indicate that global cooperation may have mobilizing potential even when it remains weakly salient in ordinary political competition \citep{fabre_public_2026}.

% AF: déjà dit plus haut, enlever
Support also persists when cooperation is embedded in a costly policy package. Overall, 68\% prefer a scenario combining global cooperation, millionaire taxation, warming limited to +2\textdegree{}C, vehicle electrification, and doubled prices for heating fuel, air travel, and beef to a status quo leading to +3\textdegree{}C warming by 2100 \citep{fabre_public_2026}. This finding demonstrates substantial acceptance of a coherent but costly package addressing climate change and global inequality, although it does not identify which component drives that support.

% AF: je reformulerais en mode: Une majorité est prête pour un gouvernement mondial, dès lors que celui-ci serait démocratique et limiterait sa compétence aux enjeux mondiaux, conformément au principe de subsidiarité. 
Taken together, these results complement the findings of \citet{ghassim_who_2024}. Public acceptance is greater when global governance is associated with clearly defined objectives, recognizable institutional reforms, and democratic procedures. The evidence does not establish unconditional demand for a centralized world government; rather, it indicates conditional acceptance of democratic, functionally limited, and policy-oriented forms of global governance \citep{fabre_public_2026,ghassim_who_2024}.

% AF: pas besoin de faire trop de paragraphes de liaison comme ça
Acceptance of global institutions in principle does not, however, establish that citizens will bear the material costs generated by the policies those institutions adopt. Climate policy provides a particularly demanding test because its private costs, distributional consequences, and dependence on participation by other countries can be made explicit.



\section{International climate policy: majority support even when people aware of the costs}

% AF: premier paragraphe pas très informatif, on ne comprend pas ce que font les papiers cités
The preceding section showed that citizens may accept supranational coordination when it is democratic, credible, and confined to genuinely global issues. A separate objection concerns the private and fiscal costs of the policies that such institutions would implement. Support for international redistribution and climate action may be overstated when respondents do not understand the implications for household income, energy prices, or domestic public budgets. The analysis must therefore move from abstract approval to the evaluation of explicitly costly policies. The relevant question is not simply whether citizens bear a cost, but how that cost is distributed, whether the policy is expected to reduce emissions effectively, and whether other countries contribute to the collective effort. Each dimension can affect acceptance, although broad international participation is not always necessary for substantial support \citep{bechtel_mass_2013,bernauer_how_2015,mcgrath_how_2017}.

Combining climate mitigation with redistribution is not merely a matter of political framing. Without compensatory transfers, ambitious mitigation policies may increase extreme poverty, whereas progressive recycling of carbon-pricing revenues and additional international climate finance can offset these regressive effects \citep{soergel_combining_2021}. The Global Climate Scheme (GCS) evaluated by \citet{fabre_public_2026} provides a demanding test of the political acceptability of such an integrated policy. Under the GCS, revenues from a global cap-and-trade system are distributed equally across adults, lifting 600 million people out of extreme poverty while imposing an average net cost of \$90 per month in the United States and \texteuro45 per month in Germany. The International Climate Scheme (ICS) applies the same principle under partial international participation. Respondents evaluate Low-, Mid-, and High-coverage scenarios accounting for 33\%, 56\%, and 72\% of global emissions, respectively, as well as a fourth variant displaying country-level distributive effects 
in the High-coverage scenario. % AF: added

% AF: il faut d'abord citer le soutien au GCS dans le papier NHB
Despite its explicit costs, the GCS is accepted by 55\% of respondents, with country-level acceptance ranging from 49\% in the United States and Russia to 85\% in Saudi Arabia \citep{fabre_public_2026}. Acceptance of the ICS increases only modestly with country coverage, from 65\% in the Low scenario to 68\% in the High scenario, a difference of 4 %3.8 AF: changed
percentage points. Among European respondents, acceptance rises when China joins the coalition but does not increase further when other high-income countries participate. By contrast, respondents in the United States and Japan react more strongly % AF: more strongly or only?
to the combined participation of China, the European Union, and other high-income countries. The effect of coverage is entirely driven by the 74\% of respondents who correctly understand that the policy would increase gasoline prices. Even under low coverage, respondents appear willing to bear the ICS's higher cost in exchange for guaranteed poverty alleviation and decarbonization in the Global South. Its greater acceptance relative to the GCS may reflect the greater credibility of a proposal that explicitly identifies participating countries, although the evidence does not isolate this mechanism.

% Overall, the evidence does not support a mechanical rule under which majority acceptance follows whenever two of three objectives---fairness, environmental effectiveness, and individual self-interest---are satisfied. These dimensions instead interact. % AF: je ne comprends pas pourquoi tu écris ça : c'est justifié par quels résultats ? Dans Douenne & Fabre (22), on trouve un soutien majoritaire dès qu'on a deux perceptions sur trois. Peut-être que d'autres trouvent autre chose, mais dans ce cas-là il faut citer et justifier
% AF: faut citer de façon plus précise les papiers ci-dessous
Higher individual costs generally reduce support, whereas equitable burden-sharing, broader participation, and credible enforcement increase the perceived legitimacy and effectiveness of international climate agreements \citep{bechtel_mass_2013}. In the specific context of North--South climate finance, acceptance is also higher when other donor countries contribute, allocation decisions % AF: comment ça allocation decisions? peu clair
are shared between donors and recipients, and recipient institutions are perceived as effective \citep{gampfer_obtaining_2014}. International reciprocity is not, however, an absolute prerequisite: substantial support for unilateral climate action persists even when other countries do not adopt equivalent policies \citep{bernauer_how_2015,tingley_conditional_2014,mcgrath_how_2017}. Perceived fairness and environmental effectiveness may therefore partly offset the influence of narrow material self-interest. The political challenge is not to identify a ``perfect'' policy, but to design a credible and transparent scheme whose environmental benefits, redistributive consequences, and burden-sharing arrangements are clearly understood.

These findings establish that costly international climate policies can attract conditional majority acceptance. The remaining question is whether this acceptance extends to a willingness to contribute a clearly specified share of personal income.


\section{Willingness to contribute: majority willing to contribute a few percent}

The preceding section showed that support for international climate policy can persist when its costs are made explicit. This section moves from acceptance of a collective policy to stated willingness to bear a personal contribution. The two concepts should not be conflated: policy acceptance measures whether respondents endorse a collective rule, whereas willingness to contribute captures the personal sacrifice they report being prepared to make. 
% The incidence of that sacrifice is equally important. % AF: deletion: unclear
% A distributionally fair scheme should protect low-income households and concentrate its net burden on those with greater ability to pay rather than treating a uniform contribution as distributionally neutral.

Using a globally representative survey of nearly 130,000 respondents across 125 countries, \citet{andre_globally_2024} ask whether individuals would contribute 1\% of their personal income to climate action. Globally, 69\% report that they would contribute, although country-level willingness ranges from 30\% to 93\%, the result is therefore majoritarian worldwide but not in every country. % AF: dire explicitement quels pays sont sous 50% et à combien ils sont (et dire qu'ils passent > 50% lorsque le coût est <1%)
The most consequential finding concerns second-order beliefs: respondents estimate that only 43\% of others would contribute, creating a 26-percentage-point gap between reported and perceived willingness. Barriers to collective climate action may therefore arise not only from private costs but also from systematic underestimation of the social support on which cooperation depends. Consistent with this interpretation, \citet{mildenberger_beliefs_2017} document widespread underestimation of pro-climate beliefs among mass publics and elites in the United States and China and show experimentally that correcting second-order beliefs increases support for climate-policy action.
% TODO: On WTP, World Bank (09) has similar results (p. 16). https://documents1.worldbank.org/curated/en/492141468160181277/pdf/520660WP0Publi1und0report101PUBLIC1.pdf (also interesting things p. 5, 10, 11, 13, 16)
% Results in Figures 1, 2 and 3 and Table S3. 11 countries don’t have majority WTC, incl. the UK (48%), the U.S. (48%), Canada (49%), Japan (49%), Israel (37%), Pakistan, New Zealand, Kazakhstan, Russia. Pluralistic ignorance is ~19 pp in U.S. and Eu4, i.e. similar to what we find for the Eu4 in Fabre, Douenne, Mattauch (but larger for the U.S.).
% “Would you be willing to contribute 1% of your household income every month to fight global warming? This would mean that you would contribute [$1] for every [$100] of this income.” Yes/No/PNR

Information does not, however, affect every form of international redistribution uniformly. In the United States, \citet{nair_misperceptions_2018} finds that information about the global income distribution increases support for foreign aid and raises real-stakes donations to foreign charities by 55\% relative to the control group. By contrast, \citet{fehr_your_2019} show that Germans systematically underestimate their position in the global income distribution and express a positive willingness to pay for information about their national and global income ranks. Experimentally correcting global-rank misperceptions nevertheless affects neither support for policies addressing global inequality nor giving to the global poor. The demand for information documented by \citet{fehr_your_2019} should therefore not be interpreted as evidence of a greater willingness to contribute. Taken together, these studies indicate that information can activate international distributive concern in some settings, but that correcting perceived global rank is not a general mechanism for increasing support for global redistribution.

Willingness to contribute remains cost-sensitive, but cost is only one determinant of policy acceptance. \citet{bechtel_mass_2013} show that higher household costs reduce support for global climate agreements, whereas equitable burden-sharing, broad participation, and credible enforcement increase it. % AF: give numbers
At the same time, support for unilateral climate action remains substantial when costs and free-riding are made explicit and appears to depend more strongly on expected effectiveness and domestic co-benefits than on strict reciprocity \citep{bernauer_how_2015,mcgrath_how_2017}. % AF: detail how they arrive at this conclusion
Support for North--South climate finance is similarly higher when recipient governments are perceived as effective, allocation decisions are jointly governed, and other donor countries contribute substantial resources \citep{gampfer_obtaining_2014}. Willingness to contribute is therefore neither unconditional nor reducible to price; it depends on policy incidence, expected effectiveness, institutional credibility, and beliefs about collective participation.

Perceptions further condition willingness to contribute internationally. Citizens tend to overestimate the foreign aid already provided by their country \citep{kaufmann_foreign_2019}, while individuals in high-income countries frequently underestimate their position in the global income distribution \citep{nair_misperceptions_2018,fehr_your_2019}. The null treatment effect documented by \citet{fehr_your_2019} nevertheless shows that correcting the latter misperception is insufficient, by itself, to increase support for global redistribution. More broadly, \citet{fabre_public_2026} finds that global climate and redistributive policies can attract majority acceptance even in countries whose citizens are expected to be net contributors. Across this literature, acceptance is stronger when burdens are distributed fairly, policies are expected to be effective, and implementation arrangements are institutionally credible \citep{fabre_public_2026,bechtel_mass_2013,gampfer_obtaining_2014}. The evidence therefore points to a conditional willingness to contribute rather than to unconditional or politically intense demand. Moreover, this willingness may remain politically latent when citizens underestimate its prevalence or attach insufficient electoral salience to international redistribution.

\section{Explaining the lack of prominence}

The preceding sections indicate that citizens are not uniformly opposed to global redistributive and climate policies. % AF: c'est bcp trop conservateur comme restitution: majorities generally accept well-designed internationally redistributive policies semble une représentation plus fidèle des résultats
On the contrary, several such policies attract conditional but frequently majority support, including when their costs and institutional design are made explicit. This creates a %central AF: delete
puzzle: why do apparently acceptable policies remain marginal in electoral competition and public debate \citep{fabre_majority_2025,fabre_public_2026}? The finding is not confined to direct survey questions. Support remains visible, although its magnitude varies, in list experiments, petition decisions, conjoint evaluations, and policy-prioritization tasks \citep{fabre_majority_2025}.

Addressing this puzzle requires distinguishing public support from political prominence and issue salience. Political prominence concerns an issue's visibility in party platforms, electoral campaigns, and public debate, whereas issue salience captures the importance citizens assign to that issue and the weight it carries in their political decisions. A policy may therefore attract majority support when evaluated in isolation without becoming sufficiently important to influence candidate evaluations or vote choice. Research on attitude importance shows that policy preferences exert greater influence on electoral behavior when citizens consider the underlying issue personally important \citep{krosnick_role_1988,dennison_review_2019}.

Beyond limited electoral salience, % AF: "limited electoral salience" is not substantiated
the gap between first-order preferences and second-order beliefs provides another potential explanation. As shown in the preceding section, citizens' reported willingness to contribute to climate action substantially exceeds their estimates of others' willingness to do so \citep{andre_globally_2024}. These misperceptions can create a coordination problem: individuals may refrain from expressing or acting on a favorable preference when they expect little participation from others. Consistent with this mechanism, \citet{mildenberger_beliefs_2017} find that members of the public and political elites systematically underestimate the prevalence of pro-climate positions and that experimentally correcting these second-order beliefs increases support for climate-policy action.

A similar mechanism may constrain political supply. \citet{fabre_public_2026} suggests that the limited supply of global redistributive proposals may reflect pluralistic ignorance among policymakers and activists, who may underestimate universalistic attitudes, public support, or the electoral returns to endorsing such policies. More general evidence shows that political representatives frequently misperceive constituency opinion and often infer that voters are more conservative than they actually are \citep{broockman_bias_2018,pilet_politicians_2023}. Although these studies do not directly examine global redistribution, they provide evidence for the plausibility of this supply-side mechanism.

% AF: ce qui serait intéressant ici c'est de comparer les expériences de donation de Nair (Respondents in the control group elected to keep 61.7% of their bonus, while donating 28.0% to a domestic charity and donating 10.1% to an international charity.) et budget allocation task de NHB et Fabre 26, ainsi que les expérience donation de Fehr et NHB (et d'autres s'il y en a d'autres). J'ai pas bcp cherché mais pas sûr que Fehr donnent leur stats descriptives (pour comparer les distributions -ou au moins les moyennes- de dons quand le destinataire est un compatriote ou un étranger): leur demander si elles ne sont pas dans le papier. 
Information alone does not necessarily transform favorable attitudes into political demand. In the United States, correcting misperceptions of relative global affluence increases support for foreign economic assistance and raises real-stakes donations to international charities, but produces only a modest increase in willingness to petition elected representatives \citep{nair_misperceptions_2018}. Respondents also continue to donate substantially more to domestic than to international charities. Similarly, although Germans systematically underestimate their position in the global income distribution, correcting this misperception affects neither support for global redistribution nor donations to the global poor \citep{fehr_your_2019}. These findings indicate that national reference groups may remain more politically influential than comparisons with the global income distribution.

Material circumstances can reinforce this national framing. Economically vulnerable left-leaning citizens may consider assistance to poorer countries morally justified while remaining reluctant to increase public expenditure when foreign aid is perceived as competing with domestic priorities \citep{cho_ideology_2024}. % AF: ça veut dire quoi "may"? on veut des résultas, des chiffres
% Misperceptions of existing aid cannot fully explain this reluctance. Citizens in donor countries substantially overestimate current aid expenditure but, on average, still prefer a level exceeding both actual and perceived aid \citep{kaufmann_foreign_2019}. % AF: c'est pour la section 2 ça

The limited prominence of international redistribution therefore cannot be reduced to widespread public opposition. % AF: c'est pas le sujet ici, cette section c'est sur pourquoi il y a une faible visibilité de la redistribution mondiale malgré son soutien majoritaire
It instead reflects the interaction of low electoral salience, pessimistic beliefs about others' support, nationally bounded reference groups, material constraints, and a political supply that may itself respond to biased perceptions and unequal political influence.

\section*{Conclusion}
This chapter began with a puzzle: % AF: non, on ne pouvait pas avoir en tête le puzzle avoir de connaître les résultats. Cf. le plan de la conclusion que j'ai proposé
internationally redistributive policies remain politically marginal even though citizens are not systematically opposed to transferring resources across borders, coordinating policy effort internationally, or delegating authority beyond the nation-state. The literature reviewed here shifts the explanation away from a simple presumption that distributive preferences are nationally bounded. Moral universalism provides a meaningful, although unevenly distributed, basis for extending distributive concern beyond compatriots \citep{enke_moral_2023}. This disposition does not translate mechanically into support for any particular policy, but it makes cross-border redistribution politically conceivable. When attention moves from abstract values to concrete instruments, the cumulative evidence supports a qualified conclusion: foreign aid, international taxation, debt relief, climate finance, and institutional cooperation can attract majority or plurality support when their objectives, incidence, and implementation are clearly specified \citep{fabre_majority_2025,fabre_public_2026}.

The successive strands of evidence point to a common interpretation. Citizens appear willing to endorse international redistribution not as an unconditional transfer of resources, but as a rule-governed collective undertaking. Support is generally stronger when beneficiaries are identifiable, resources are protected against diversion, burdens fall on those with greater capacity to contribute, and international participation is reciprocal and credible \citep{fabre_majority_2025,fabre_public_2026,gampfer_obtaining_2014}. The same logic extends from fiscal instruments to political institutions. Unspecified transfers of authority beyond the nation-state generate resistance, whereas democratic and functionally limited forms of global governance attract substantially greater acceptance \citep{ghassim_who_2024}. Similarly, making private costs explicit reduces support but does not eliminate it when policies combine environmental effectiveness, progressive incidence, and visible poverty reduction \citep{bechtel_mass_2013,fabre_public_2026}. Citizens therefore appear to evaluate not international redistribution in isolation, but the institutional package through which it is financed, allocated, implemented, enforced, and justified.

Returning to the contrast introduced at the outset, the moral and geographical distance separating Clara and Moussa does not constitute an absolute barrier to redistribution. The decisive question is whether concern for Moussa can be translated into a policy that Clara perceives as fairly financed, effectively implemented, and supported by legitimate and credible institutions.

The theoretical implication is that preferences over international redistribution cannot be adequately represented along a single cosmopolitan--nationalist dimension. They are multidimensional and conditional on interactions among moral concern, material incidence, expected effectiveness, institutional legitimacy, and beliefs about collective participation. Fairness operates vertically through the distribution of costs across income groups and horizontally through the allocation of responsibilities across countries. Effectiveness matters because citizens distinguish transfers that plausibly reach their intended beneficiaries from those exposed to corruption, free-riding, or weak enforcement. Institutional credibility connects these dimensions: democratic procedures, transparent revenue allocation, and identifiable participation can help translate a general willingness to assist distant beneficiaries into acceptance of a specific policy. This framework also clarifies why self-interest matters without being determinative. Higher private costs and competition with domestic spending constrain support, yet they coexist with substantial willingness to contribute when policies are perceived as equitable, effective, and institutionally credible.

Important tensions nevertheless remain. Much of the evidence relies on stated preferences elicited after respondents receive information and evaluate policy designs that are more explicit than those ordinarily encountered in electoral debate. Majority acceptance in surveys therefore need not imply intense preferences, sustained attention, or a willingness to sanction parties that neglect these issues. The findings also reveal substantial heterogeneity across countries, ideological groups, and income positions, while information treatments produce context-dependent effects. Correcting misperceptions increases donations or support for foreign aid in some settings, but information about citizens' global income rank does not consistently alter redistributive preferences \citep{nair_misperceptions_2018,fehr_your_2019}. Moreover, willingness to contribute may remain politically latent when individuals underestimate its prevalence, as the gap between first- and second-order beliefs in climate cooperation illustrates \citep{andre_globally_2024,mildenberger_beliefs_2017}. These limitations do not overturn the chapter's central conclusion; rather, they identify mechanisms through which favorable attitudes may fail to become organized political demand.

The central unresolved question is therefore less whether a constituency for international redistribution exists than under what conditions its latent and conditional acceptance becomes politically consequential. Future research should examine how issue salience is created, how citizens update beliefs about others' preferences, how parties and representatives infer electoral demand, and which institutional features sustain support outside survey settings. It should also distinguish more systematically among policy approval, willingness to pay, budgetary prioritization, political participation, and vote choice, each of which captures a different stage in the translation of moral concern into public policy. This agenda redirects attention from the measurement of isolated attitudes toward the political processes connecting public preferences, institutional design, and policy supply. Explaining this translation---and the conditions under which it succeeds or fails---is the next step toward understanding why internationally redistributive policies remain scarce despite a broader social basis of acceptance than prevailing political debates suggest.


\newpage
\bibliographystyle{plainnat}
\bibliography{references}

\end{document}